% GENERATED FILE. Edit canonical lessons, then run npm run generate:books.
% canonical-source-hash: fnv1a64:5a5a45557cc23b47
% canonical-lessons: PA-C42-oh, PA-C42-asi, PA-C42-oh-write, PA-C42-jo, PA-C42-jo-oh

\chapter{He, We, and the One Who}
\label{ch:pa-pronouns-and-relative}
\hlchaptermodality{42}

\begin{chapteropening}
\textbf{By the end of this chapter:} \emph{I can talk about somebody who is not present, say we, and describe a person instead of naming them.}

The absent third person, the missing we, and the relative. asi is the oddest gap in the inventory: the corpus teaches the first-person PLURAL future ending four separate times -- milaange, phir milaange, jaldi milaange, kall milaange -- so a reader has been saying we will meet for chapters without ever owning the word for we. jo completes the j- family and pairs with oh in the same two-ended shape as je ... taan.
\end{chapteropening}

\section[oh]{\pa{ਉਹ} (oh) — he, she, it}

\label{lesson:PA-C42-oh}



\begin{quote}
From the chapter behind you, the word that names a purpose. And from chapters 2 and 3, the words for \textbf{I} and for \textbf{you}, both of them.

\begin{itemize}
\raggedright
  \item \emph{Recall:} \emph{laī} — then \emph{maiṁ}, and \emph{tū̃} / \emph{tusī̃}
  \item \emph{Recall:} \emph{ki} — from two chapters back
  \item \emph{Recall:} \emph{nahī̃} before the verb — from five chapters back, cold
\end{itemize}

You have been able to talk about yourself and about the person in front of you since you first said \textbf{\pa{ਮੈਂ}}. You have never been able to talk about anybody \textbf{absent}.
\end{quote}

\begin{cousinweb}
\begin{quote}
\textbf{\pa{ਉਹ}} — \emph{oh} — \textbf{he}, \textbf{she}, or \textbf{it}
\end{quote}

One word, three English ones. Punjabi's third person does not make you choose a gender before you speak.

\textbf{\pa{ਉਹ}} also means \emph{that}, pointing at something away from both of you. The pronoun and the pointing word are one word — which is the same doubling you have already seen, where \textbf{\pa{ਤਾਂ}} at the head of a result was a demonstrative doing a joining job.

With this the pronoun set is complete: \textbf{\pa{ਮੈਂ}}, \textbf{\pa{ਤੂੰ}} / \textbf{\pa{ਤੁਸੀਂ}}, \textbf{\pa{ਉਹ}}.
\end{cousinweb}

\subsection*{Your turn}
\begin{itemize}
\raggedright
  \item \emph{Say it:} \emph{maiṁ}, \emph{tusī̃}, \emph{oh} — the three in a row.
  \item \emph{Build:} a sentence about someone who is not in the room.
  \item \emph{Say it:} \emph{oh} as \emph{that}, pointing across the room.
\end{itemize}

\begin{tcolorbox}[breakable,colback=teal!4,colframe=teal!35!black,title={Before you move on}]
How many English pronouns does \textbf{\pa{ਉਹ}} cover? (Three.) What else does the same word do? (Points: \emph{that}.)
\end{tcolorbox}
\section[asīṁ]{\pa{ਅਸੀਂ} (asīṁ) — we}

\label{lesson:PA-C42-asi}



\begin{quote}
One lesson back, the pronoun for someone absent. And from the farewell chapters: say \emph{we will meet soon}, and \emph{we will meet tomorrow}.

\begin{itemize}
\raggedright
  \item \emph{Recall:} \emph{oh} — then \emph{jaldī milāṁge}, \emph{kall milāṁge}
  \item \emph{Recall:} \emph{maiṁ sochdā hāṁ ki} — from two chapters back
  \item \emph{Recall:} \emph{mainū̃ samajh nahī̃ āī} — from five chapters back, cold
\end{itemize}
\end{quote}

\begin{cousinweb}
\begin{quote}
\textbf{\pa{ਅਸੀਂ}} — \emph{asīṁ} — \textbf{we}
\end{quote}

Here is the odd thing, and it is worth stopping on. You have been saying \textbf{we} for chapters. The \textbf{-\pa{ਾਂਗੇ}} ending on \emph{milāṁge} is the first-person \textbf{plural} future — \emph{we will meet} — and this book has taught it four separate times.

The ending was there. The pronoun was not. So a reader has been saying \emph{we will meet} over and over without ever owning the word for \emph{we}.

Punjabi often leaves the pronoun out when the ending already carries it, which is why nothing sounded wrong. But you cannot say \textbf{we} on its own, or put it in front of anything else, until you have the word.
\end{cousinweb}

\subsection*{Your turn}
\begin{itemize}
\raggedright
  \item \emph{Say it:} \emph{asīṁ}, alone.
  \item \emph{Say it:} \emph{asīṁ … milāṁge} — the pronoun and the ending together.
  \item \emph{Contrast:} \emph{maiṁ} and \emph{asīṁ} — one, then more than one.
\end{itemize}

\begin{tcolorbox}[breakable,colback=teal!4,colframe=teal!35!black,title={Before you move on}]
Say \emph{we}. Which ending have you been using for chapters that already meant it? (The \textbf{-\pa{ਾਂਗੇ}} of \emph{milāṁge}.)
\end{tcolorbox}
\section[oh]{\pa{ਉਹ} reaches the page}

\label{lesson:PA-C42-oh-write}



\begin{quote}
One chapter back you wrote \textbf{\pa{ਜੇ}}. Write it once from memory, then say the word for \emph{we}.
\end{quote}

\subsection*{Script — two signs, neither new}
\begin{quote}
\textbf{\pa{ਉ}} — the \emph{u} carrier · \textbf{\pa{ਹ}} — \emph{ha}
\end{quote}

Both taught long ago. \textbf{\pa{ਉ}} is the carrier you met when the book first showed a vowel standing on its own, and \textbf{\pa{ਹ}} you have written since \textbf{\pa{ਨਹੀਂ}} at the top of this run.

\subsection*{Writing — guided copy, model in view}
Copy \textbf{\pa{ਉਹ}} once with the model in view. Compare its first letter against the \textbf{\pa{ਅ}} of \textbf{\pa{ਅਤੇ}} — two different carriers, and the difference matters.

\begin{tcolorbox}[breakable,colback=teal!4,colframe=teal!35!black,title={Before you move on}]
Write \textbf{\pa{ਉਹ}} from memory. How many of its signs were new? (None.)

\begin{itemize}
\raggedright
  \item \emph{Recall:} \textbf{\pa{ਕਿ}} written — from two chapters back
  \item \emph{Recall:} \textbf{\pa{ਨਹੀਂ}} written — from five chapters back, cold
\end{itemize}
\end{tcolorbox}
\section[jo]{\pa{ਜੋ} (jo) — the one who}

\label{lesson:PA-C42-jo}



\begin{quote}
One lesson back you wrote \textbf{\pa{ਉਹ}}. One chapter back, the two words that open a clause with \textbf{\pa{ਜ}}.

\begin{itemize}
\raggedright
  \item \emph{Recall:} \textbf{\pa{ਉਹ}}, then \emph{je}, \emph{jadoṁ}
  \item \emph{Recall:} \emph{kyoṁ} — from two chapters back
  \item \emph{Recall:} \emph{mainū̃ nahī̃ patā} — from five chapters back, cold
\end{itemize}
\end{quote}

\begin{cousinweb}
\begin{quote}
\textbf{\pa{ਜੋ}} — \emph{jo} — \textbf{the one who}, \textbf{the thing that}
\end{quote}

The third member of the family. \textbf{\pa{ਜੇ}} supposes, \textbf{\pa{ਜਦੋਂ}} times, and \textbf{\pa{ਜੋ}} describes — and all three open their clause with the same letter.

\begin{quote}
\textbf{\pa{ਜੋ}} \emph{(what they do)} … — \textquotedblleft{}the one who …\textquotedblright{}
\end{quote}

Punjabi marks a describing clause at its front, the way it marked a condition with \textbf{\pa{ਜੇ}}. Once you hear the \textbf{\pa{ਜ}}, you know a clause is starting and that something will come back to be pointed at.
\end{cousinweb}

\subsection*{Your turn}
\begin{itemize}
\raggedright
  \item \emph{Say it:} \emph{jo}, then something a person might do.
  \item \emph{Run through:} the three \textbf{\pa{ਜ}} words in a row — \emph{je}, \emph{jadoṁ}, \emph{jo}.
  \item \emph{Build:} describe someone instead of naming them.
\end{itemize}

\begin{tcolorbox}[breakable,colback=teal!4,colframe=teal!35!black,title={Before you move on}]
Describe a person rather than naming them. What do the three \textbf{\pa{ਜ}} words have in common? (Each opens its own clause.)
\end{tcolorbox}
\section[jo … oh …]{\pa{ਜੋ} … \pa{ਉਹ} … (jo … oh …) — the one who … that one …}

\label{lesson:PA-C42-jo-oh}



\begin{quote}
This chapter's three words, and the two-part condition from the chapter before.

\begin{itemize}
\raggedright
  \item \emph{Recall:} \emph{oh}, \emph{asīṁ}, \emph{jo} — and \emph{je … tāṁ …}
  \item \emph{Recall:} \emph{kyoṁki} — from two chapters back
  \item \emph{Recall:} the repair kit — from five chapters back, cold
\end{itemize}
\end{quote}

\begin{cousinweb}
\begin{quote}
\textbf{\pa{ਜੋ}} \emph{(description)}, \textbf{\pa{ਉਹ}} \emph{(the one meant)}
\end{quote}

This is the same shape as \textbf{\pa{ਜੇ} … \pa{ਤਾਂ} …}, and seeing that is most of the work. Punjabi marks \textbf{both} ends: the clause that describes, and the thing described. English collapses them — \emph{the one who came is my friend} — and Punjabi keeps them visibly paired.

Now the whole system is visible:

\begin{quote}
\textbf{\pa{ਜ}-} opens the clause · \textbf{\pa{ਉ}-} / \textbf{\pa{ਤ}-} points at what it refers to
\end{quote}

\textbf{\pa{ਜੇ} … \pa{ਤਾਂ}}, \textbf{\pa{ਜੋ} … \pa{ਉਹ}}. One habit, and you have met it twice.
\end{cousinweb}

\subsection*{Your turn}
\begin{itemize}
\raggedright
  \item \emph{Build:} \emph{jo} + a description, \emph{oh} + who is meant.
  \item \emph{Contrast:} the same idea with a plain name, then with the pair.
  \item \emph{Say it:} \emph{je … tāṁ …} then \emph{jo … oh …} — hear that they are one shape.
\end{itemize}

\begin{tcolorbox}[breakable,colback=teal!4,colframe=teal!35!black,title={Before you move on}]
Describe someone with the pair. Which earlier pattern is this the same shape as? (\textbf{\pa{ਜੇ} … \pa{ਤਾਂ} …}.) Next chapter: the question words, all of them built on one letter.
\end{tcolorbox}
